% ==== P4 draft F1 — §15.1 대체 블록 ====
\subsection{왜 흑연에는 없고 LCO 에는 있는가 --- MIT 의 물리}\label{sec:lco-why}
삽입 반응의 전체 엔트로피 변화는 ``리튬 자리 배열''(config)$\cdot$``격자 진동''(vib)$\cdot$``전자 준위 점유''
(electronic) 세 자유도의 합이며, 흑연에서는 셋째 몫이 작다 --- 흑연$\cdot$$\mathrm{LiC_6}$ 모두 전도성이 좋아
충방전 동안 $g(E_F)$ 가 크게 바뀌지 않아 전자 분포의 엔트로피 \emph{변화}가 거의 없다. LCO 는 다르다:
$x{=}1$($\mathrm{LiCoO_2}$)은 Co$^{3+}$($t_{2g}^6$ low-spin, 닫힌 껍질)라 \emph{전기 절연체}
($g(E_F)\!\approx\!0$)이고, 그 물리 부류는 산소 팔면체의 \emph{결정장}(crystal field)이 가른 아래 $t_{2g}$
띠를 $d^6$ 전자 여섯이 남김없이 채워 가득 찬 띠와 빈 $e_g$ 띠 사이 갭만 남는 \emph{밴드 절연체}(band
insulator)다. 탈리튬화로 $x$ 가 $\sim$0.94 아래로 내려가면 Co$^{4+}$ 정공이 생겨 $t_{2g}$ 띠에 전도 전자가
열려 \emph{금속}이 된다($g(E_F)\!\to\!$유한)\cite{menetrier1999,motohashi2009}.

이 MIT 는 $x\approx0.75$--$0.94$ 의 2상역에서 일어나며(\S\ref{sec:lco-hys} 의 T1) 그 구간에서만 전자 분포가
켜지므로 전자 엔트로피 \emph{변화}도 그 구간에 국소한다. 다만 ``왜 첫 정공에서가 아니라 $\sim$0.94
부터인가''는 단순 띠 채움 논리 바깥의 질문이다 --- 그 논리라면 가득 찬 띠에 정공이 조금이라도 생기는 즉시
금속이어야 하는데, 실측은 절연상$\cdot$금속상이 공존하는 2상역을 거쳐 넘어가는 \emph{1차 전이}이기
때문이다\cite{menetrier1999,reimers1992}. 이 물음의 해소로 제시된 미시 기작(Li 빈자리에 속박된 정공 불순물
준위의 Mott 전이\cite{marianetti2004})까지 포함해, 금속--절연체 전이의 물리 배경은 별도의 배경 박스가
자족적으로 정리한다.

\begin{bgbox}
\textbf{금속--절연체 전이와 Mott 물리 ---} 고체가 절연체가 되는 길은 두 부류로 갈린다 ---
\emph{밴드 절연체}(band insulator)는 전자가 띠를 남김없이 가득 채우고 다음 빈 띠까지 갭이 열려 있어,
전자 사이 상호작용을 묻지 않는 단일입자 띠 그림 안에서 이미 전류가 막히는 절연체다.
\emph{Mott 절연체}(Mott insulator)는 그 그림이 틀리는 절연체다 --- 띠가 부분만 차서 띠 그림으로는
금속이어야 하나, 한 자리를 두 전자가 겹쳐 점유할 때 드는 전자간 Coulomb 반발 $U$ 가 이웃 자리로 건너뛰는
\emph{이동적분}(hopping integral) $t$ 를 누르면 전자가 자리마다 하나씩 국재해 절연체가
된다\cite{mott1968,imada1998}(이것이 전자 \emph{상관}(correlation) 효과이며, 여기의 $U$ 와 $t$ 는 본 문건의
개회로전압 $U_j$ 와 별개인 상관 물리의 관용 기호다). 금속--절연체 전이(MIT)는 이 두 부류 사이 경계를 넘는
사건이며, 띠의 채움(도핑)$\cdot$띠폭(압력)$\cdot$상관의 경쟁으로 일어난다\cite{imada1998}.
$\mathrm{Li}_x\mathrm{CoO_2}$ 의 $x{=}1$ 끝점($\mathrm{LiCoO_2}$, Co$^{3+}$ $t_{2g}^6$ low-spin)은 가득 찬
$t_{2g}$ 띠가 갭 아래 닫힌 밴드 절연체다 --- 절연의 근거가 상관이 아니라 띠 채움이라, 반쯤 찬 띠가 $U$ 로
얼어붙는 Mott 절연체와는 부류가 다르다\cite{menetrier1999}. 탈리튬화 정공은 이 구분을 시험대에 올린다 ---
띠 그림은 가득 찬 띠에서 전자를 덜어내는 즉시 금속을 예상하나, 실측 MIT 는 $x\approx0.75$--$0.94$ 의
2상역에서 절연상$\cdot$금속상이 공존하는 \emph{1차 전이}다\cite{menetrier1999,reimers1992,motohashi2009}.
Marianetti--Kotliar--Ceder 는 대형 초격자 제일원리 계산으로 이 어긋남의 해소를 제시했다: 묽은 정공은 host
$t_{2g}$ 띠를 떠돌지 않고 Li 빈자리의 인력에 속박된 \emph{불순물 준위}(impurity state)를 이루며, 그
준위들이 겹친 \emph{불순물 띠}(impurity band)가 임계 농도에서 비국재화하는 Mott 전이가 관측된 1차
MIT$\cdot$2상 공존을 설명한다는 것이다\cite{marianetti2004}. 곧 상관 물리는 host 띠가 아니라 불순물
준위에서 일하며, ``$x{=}1$ 끝점은 밴드 절연체''라는 부류 판정과 ``전이는 Mott 형 1차 전이''라는 실측이
이렇게 양립한다. 본 문건의 forward 모델은 이 미시 기작 전체를 조성축 게이트 $g(E_F,x)$ 하나로 현상학화한다
--- 게이트 중심 $x_\mathrm{MIT}$ 는 실측 2상역의 중앙에, 폭 $\Delta x_\mathrm{MIT}$ 는 결함이 퍼뜨리는 발현
조성의 분산에 대응한다(수치 초기값과 logistic 형태의 정당화는 \S\ref{sec:lco-gate} 가 닫는다). 미시 기작은
게이트가 \emph{왜} 그 조성에서 켜지는가에 답하고, 게이트는 그 켜짐이 forward 모델에 \emph{어떻게}
들어가는가에 답한다.
\end{bgbox}

곧 ``전자항이 켜지는 게이트''는 MIT 의 직접 결과이며, 그 게이트에 실리는 전자 엔트로피의 함수형은 다음
소절(\S\ref{sec:lco-Se})이 Fermi--Dirac 분포에서 유도한다.
% ==== 블록 끝 ====
% NOTE (F1 자기보고):
% [Read Coverage] AUTHOR_BRIEF.md 전문 / ch1_sec15_lcoelec.tex 전문(§15.1 L12-20, §15.4 게이트 정당화
%   L146-193 — 역할 분담 준수: bgbox 는 "왜(기작)" 만 다루고 게이트 파라미터의 수치 근거 (i)-(iv)는
%   §15.4 에 남겨, 접속 문장 1개(정당화는 \S\ref{sec:lco-gate} 가 닫는다)로만 연결 — 중복 없음) /
%   ch1_sec13_lcohys.tex L120-159(T1 슬롯 분리 L134-156 — bgbox 는 히스 gap 슬롯(Omega_1)을 건드리지 않고
%   전자 배경만 다뤄 모순 없음; L136-137 의 기작 명명 어구를 본문 handoff 가 동일하게 mirror) /
%   V1020_STYLE_RUBRIC.md 전문 / ch1_bib.tex L19-25(허용 6키 실재 확인) / ch1_sec02a_part0.tex L160-188
%   (bgbox 관용 exemplar — 제목 \textbf ... --- 형식·cite 후 괄호 부연·자족 마무리 패턴 준거) /
%   ch1_sec11_lcointro.tex L36-55(MIT 영문 병기 기출(L42) 확인 — 본문 재병기 생략, bgbox 내부는 자족
%   요건으로 "금속--절연체 전이(MIT)" 한글 anchor 만 재제시, 영문 재확장은 이중 확장 충돌 회피 위해 생략).
% [보존 확인 — 기존 §15.1 물리 문장·인용 전량 유지]
%   (1) 세 자유도(config/vib/electronic) 합 + 흑연 셋째 몫 작음(흑연·LiC6 전도성, g(E_F) 불변, 변화 거의
%       없음) — 1문단 첫 문장, 원문 그대로.
%   (2) x=1(LiCoO2) = Co3+(t2g6 low-spin, 닫힌 껍질)·전기 절연체(g(E_F)≈0) — 1문단, 원문 구절 보존 위에
%       밴드 절연체 부류 정밀화(브리프 지시 1 이행).
%   (3) 탈리튬화 ~0.94 아래 Co4+ 정공·t2g 띠 전도 전자·금속(g(E_F)→유한)\cite{menetrier1999,motohashi2009}
%       — 1문단 끝, 원문 그대로(기존 인용쌍 원위치 유지).
%   (4) MIT x≈0.75--0.94 2상역(\S\ref{sec:lco-hys} 의 T1)·전자 분포 켜짐·엔트로피 변화 국소 — 2문단 첫
%       문장, 원문 그대로(게이트 국소화 내용 보존).
%   (5) "전자항이 켜지는 게이트는 MIT 의 직접 결과" — 마무리 다리에 어구·취지 보존 + §15.2 전달
%       (브리프 지시 3 이행).
% [자체검수 — 물리 가드 3항 점검]
%   가드1 (x=1 을 Mott 절연체로 서술 금지): 본문·bgbox 모두 "밴드 절연체" 로 부류 명시. bgbox 에 허용
%     대조 가드 형식("절연의 근거가 상관이 아니라 띠 채움이라, ... Mott 절연체와는 부류가 다르다") 적용.
%     Mott 물리의 소관은 불순물 준위(불순물 띠)로 한정 서술 — 위반 없음.
%   가드2 (Marianetti 과대해석 금지): "해소를 제시했다 / 설명한다는 것이다 / 해소로 제시된" 수준의 귀속
%     서술로 한정, "확정된 유일 기작" 류 단정 없음 — 위반 없음.
%   가드3 (기존 문장·인용 유실 금지 + 형식 제약): menetrier1999·motohashi2009 기존 인용 원위치 유지(보존
%     (3)); 사용 인용은 허용 6키(mott1968·imada1998·marianetti2004·menetrier1999·reimers1992·motohashi2009)
%     뿐; 수식·라벨 신설 0(산문+bgbox 만, 기존 \label{sec:lco-why} 유지); 신규 기호 U·t 는 bgbox 내부
%     한정 + 개회로전압 U_j 와의 기호 충돌 가드 병기 — 위반 없음.
% [분량] 본문 소절 16행(문단 2개 + 마무리 다리 1문장) / bgbox 20행·9문장(요구 8-12문장·12-20행 내외 충족).
